% ATR Paper — INTERSPEECH 2027 target
% Authors: Ifeanyi Dike, Happy Nkanta Monday
\documentclass[a4paper]{article}

\usepackage{INTERSPEECH2025}   % replace with correct template when submitting
\usepackage{amsmath}
\usepackage{booktabs}
\usepackage{multirow}
\usepackage{graphicx}
\usepackage{url}
\usepackage{hyperref}

\title{ATR: An Acoustic Turn-over Rate Metric for Evaluating\\
       Overlap Handling in Full-Duplex Spoken Dialogue Systems}

\name{Ifeanyi Dike$^1$, Happy Nkanta Monday$^2$}

\address{
  $^1$Datastrut AI \\
  $^2$[Happy's affiliation]
}

\email{ifeanyi@datastrut.ai, [Happy's email]}

\begin{document}

\maketitle

%=============================================================================
\begin{abstract}
Turn-over Rate (TOR), introduced in Full-Duplex-Bench, measures how often a
full-duplex spoken dialogue model yields the speaking floor during overlapping
speech. TOR is computed from ASR transcripts, and this dependency turns out to
matter: two widely-used ASR backends produce TOR estimates that diverge by up
to 23.5 percentage points on the same audio, making cross-study comparisons
unreliable. We propose the Acoustic Turn-over Rate (ATR), which replaces the
ASR step with voice activity detection. ATR asks the same question --- did the
model's audio fall silent before the overlap window closed? --- but answers it
directly from the waveform, without transcription. Evaluated on Moshi across
all four Full-Duplex-Bench v1.5 scenarios, three independent VAD backends
(Silero, WebRTC, pyannote) agree within 7.6 percentage points on
\emph{user\_interruption}, roughly three times more consistent than ASR-based
TOR on the same audio. ATR requires no language model, runs locally, and is
language-agnostic, making it a practical and more reproducible alternative for
benchmarking full-duplex speech systems.
\end{abstract}

\noindent\textbf{Index Terms:} full-duplex dialogue, turn-taking evaluation,
voice activity detection, spoken dialogue systems, overlap handling

%=============================================================================
\section{Introduction}
\label{sec:intro}

Voice has always been the most natural interface for conversation. The appeal
of full-duplex spoken dialogue systems is obvious --- they listen and speak
simultaneously, the way people do. But this also means they have to decide, in
real time, whether to keep talking when someone else starts. That decision is
the hard part, and measuring how well a system makes it turns out to be
harder than it looks.

Full-Duplex-Bench~\cite{lin2025fullduplex} proposed TOR --- Turn-over Rate ---
as a way to quantify floor-yielding behaviour. The metric is straightforward:
for each test clip, did the model stop speaking before the user's overlap
window closed? Average across clips and you get a number between 0 and 1.
Clean idea. The issue is that figuring out ``when the model stopped speaking''
is not as simple as it sounds. TOR answers this by running ASR on the model's
audio output and taking the timestamp of the last recognised word --- which
means the evaluation result depends on which ASR system you picked. And
systems disagree, sometimes by a lot.

We ran two standard ASR backends on the same Moshi~\cite{defossez2024moshi}
outputs across all four Full-Duplex-Bench v1.5~\cite{lin2025fullduplexv15}
scenarios. On \emph{user\_backchannel}, MLX Whisper gave TOR~$= 0.44$;
AssemblyAI gave $0.20$. Same audio, 23.5 percentage point gap. On
\emph{talking\_to\_other}, the difference was 22 points. This is not a corner
case --- it is what happens every time someone replicate an experiment with a
different ASR system.

The fix we propose is to skip ASR altogether. ATR --- Acoustic Turn-over Rate
--- asks the same question as TOR but uses voice activity detection instead of
transcription. When did the model's audio go silent? VAD answers that without
needing a vocabulary, a language model, or a transcription server. We tested
three VAD backends on the same clips: Silero~\cite{silero2021},
WebRTC~\cite{webrtcvad}, and pyannote~\cite{plaquet2023,bredin2020}. They
disagreed by at most 7.6 percentage points on \emph{user\_interruption}, and
14.7 points in the worst case across all tasks --- two to three times more
stable than ASR, consistently.

%=============================================================================
\section{Background}
\label{sec:background}

\subsection{Turn-taking in spoken dialogue}

Human conversation is not a strict alternation of speaker turns. Overlapping
speech --- interruptions, backchannels, simultaneous starts --- accounts for
roughly 10 to 20\% of spoken time~\cite{defossez2024moshi}, and managing it
gracefully is one of the things that makes a conversation feel natural rather
than mechanical~\cite{sacks1974,skantze2021}. Early voice assistants sidestepped
the problem with push-to-talk and endpointing; full-duplex systems attempt to
handle it continuously.

\subsection{Full-Duplex-Bench}

Full-Duplex-Bench~\cite{lin2025fullduplex} (FDB) is a systematic evaluation
framework for full-duplex spoken dialogue models, assessing four overlap
scenarios: \emph{user interruption} (the user takes the floor while the model
is speaking), \emph{user backchannel} (the user produces a short acknowledgement
without claiming the floor), \emph{talking to others} (the user addresses
a third party), and \emph{background speech} (ambient speech not directed at
the model). For each scenario, the benchmark supplies audio stimuli and
reference timestamps marking the start and end of the overlap window.

FDB v1.0~\cite{lin2025fullduplex} computes TOR from ASR transcripts generated
by NVIDIA Parakeet-TDT-0.6B-v2. FDB v1.5~\cite{lin2025fullduplexv15}
(accepted at ICASSP 2026) drops TOR entirely and replaces it with stop and
response latency alongside categorical behaviour labels assigned by GPT-4o.
The shift away from TOR in v1.5 is suggestive; our experiments offer an
empirical explanation for why.

\subsection{ASR variance in spoken language evaluation}

The sensitivity of downstream metrics to ASR system choice is not a new
observation. Word error rates differ substantially across systems on the same
audio~\cite{radford2023whisper}, and word-level timestamp accuracy varies
further. When a metric is derived from ASR timestamps, it inherits this
variance. What has not been quantified before is how large this effect is
specifically for overlap-handling evaluation --- and whether a transcript-free
alternative can do better.

%=============================================================================
\section{TOR and Its ASR Dependence}
\label{sec:tor}

\subsection{Definition}

Following~\cite{lin2025fullduplex}, TOR is:
\begin{equation}
    \text{TOR} = \frac{1}{N} \sum_{i=1}^{N}
    \mathbf{1}\!\left[t^{\text{asr}}_i \leq t^{\text{ref}}_i\right]
    \label{eq:tor}
\end{equation}
where $t^{\text{asr}}_i$ is the end timestamp of the last ASR-recognised word
in sample $i$'s model output, and $t^{\text{ref}}_i$ is the end of the overlap
window from the benchmark metadata. TOR~$= 1$ means the model always yielded
the floor before the window closed; TOR~$= 0$ means it never did.

\subsection{Sensitivity to ASR choice}

We computed TOR for Moshi on Full-Duplex-Bench v1.5 using two backends:
\begin{itemize}
    \item \textbf{MLX Whisper} --- \texttt{mlx-community/whisper-large-v3-mlx},
          a Metal-accelerated port of Whisper large-v3~\cite{radford2023whisper}
          running locally on Apple Silicon.
    \item \textbf{AssemblyAI} --- Universal-1~\cite{assemblyai2024}, a
          commercial ASR API with word-level timestamp support.
\end{itemize}

\begin{table}[t]
  \centering
  \caption{TOR from two ASR backends on Moshi outputs. $\Delta$ is the
           absolute difference. $n$ is the number of samples with a
           non-empty transcript.}
  \label{tab:asr_sensitivity}
  \begin{tabular}{lcccc}
    \toprule
    Scenario & $n_{\text{MLX}}$ & MLX & AAI & $\Delta$ \\
    \midrule
    user\_interruption   & 196 & 0.000 & 0.000 & 0.000 \\
    user\_backchannel    &  98 & 0.439 & 0.204 & \textbf{0.235} \\
    talking\_to\_other   & 100 & 0.360 & 0.140 & \textbf{0.220} \\
    background\_speech   & 100 & 0.310 & 0.140 & \textbf{0.170} \\
    \bottomrule
  \end{tabular}
\end{table}

Results are in Table~\ref{tab:asr_sensitivity}. Three of four scenarios show
17--23.5 percentage point gaps between the two backends on the same audio.
The \emph{user\_interruption} case looks stable at first glance --- both
systems return TOR~$= 0$ --- but only because both agree on a floor value.
AssemblyAI returned an empty transcript for 108 of 200
\emph{user\_interruption} samples (54\%), because those clips contain
near-silence after the interruption endpoint and the API produces no output.
A different system might handle near-silent audio differently, pushing that
zero in either direction.

The core issue is structural. ASR systems apply different confidence thresholds,
post-processing pipelines, and endpointing heuristics. A word near the end of
an utterance that one system transcribes, another may discard. That one word
can shift $t^{\text{asr}}_i$ enough to flip the binary indicator in
Equation~\ref{eq:tor}, and those flips accumulate into the large gaps
observed in Table~\ref{tab:asr_sensitivity}.

%=============================================================================
\section{Acoustic Turn-over Rate (ATR)}
\label{sec:atr}

\subsection{Definition}

ATR replaces $t^{\text{asr}}_i$ in Equation~\ref{eq:tor} with a VAD-based
estimate of the last speech end:
\begin{equation}
    \text{ATR} = \frac{1}{N} \sum_{i=1}^{N}
    \mathbf{1}\!\left[t^{\text{vad}}_i \leq t^{\text{ref}}_i\right]
    \label{eq:atr}
\end{equation}
where $t^{\text{vad}}_i$ is the end of the last VAD-detected speech segment
in sample $i$'s model output. Everything else --- the reference timestamps,
the binary indicator, the average --- stays the same as TOR.

\subsection{VAD post-processing}

For each model output audio file we:
\begin{enumerate}
    \item Resample to 16~kHz (required by all three VAD backends).
    \item Run the VAD to get a list of $(t_{\text{start}}, t_{\text{end}})$
          speech segment pairs.
    \item Merge adjacent segments separated by $\leq 300$~ms.
    \item Set $t^{\text{vad}}_i$ to the end of the last merged segment.
          Samples with no detected speech are excluded from the ATR computation
          and counted separately as a speech-absence rate.
\end{enumerate}

The 300~ms merge threshold follows standard VAD post-processing
practice and is the same across all three backends for comparability.

%=============================================================================
\section{Experiments}
\label{sec:experiments}

\subsection{Setup}

\paragraph{Benchmark.}
We use Full-Duplex-Bench v1.5~\cite{lin2025fullduplexv15}, which provides
498 audio stimuli across four scenarios: \emph{user\_interruption} (200),
\emph{user\_backchannel} (98), \emph{talking\_to\_other} (100), and
\emph{background\_speech} (100).

\paragraph{Model.}
We evaluate \textbf{Moshi}~\cite{defossez2024moshi}
(\texttt{kyutai/moshiko-mlx-bf16}), a real-time full-duplex speech model
built on a multi-stream language model with a neural audio codec.
We run inference offline using the \texttt{moshi\_mlx} Python API and
\texttt{rustymimi} StreamTokenizer, bypassing the WebSocket server entirely.
This produces 498 audio output files, all confirmed non-silent.

\paragraph{VAD backends.}
\begin{itemize}
    \item \textbf{Silero VAD}~\cite{silero2021} --- a lightweight LSTM-based
          VAD pre-trained on multilingual speech data.
    \item \textbf{WebRTC VAD}~\cite{webrtcvad} --- a GMM-based VAD from the
          WebRTC open-source project, using 30~ms frames at aggressiveness
          level~2.
    \item \textbf{pyannote segmentation-3.0}~\cite{plaquet2023,bredin2020}
          --- a transformer-based segmentation model wrapped in a
          \texttt{VoiceActivityDetection} pipeline with default
          \texttt{min\_duration\_on} and \texttt{min\_duration\_off} set
          to 0.
\end{itemize}

All three backends receive 16~kHz mono audio and use the same 300~ms
merge threshold described in Section~\ref{sec:atr}.

\subsection{ATR stability across VAD backends}

\begin{table}[t]
  \centering
  \caption{ATR from three VAD backends on Moshi outputs. $n$ is the number
           of samples with detected speech (varies by backend). Max~$\Delta$
           is the largest pairwise absolute difference across the three backends.}
  \label{tab:vad_variance}
  \begin{tabular}{lccccc}
    \toprule
    Scenario & $n$ & Silero & WebRTC & pyannote & Max~$\Delta$ \\
    \midrule
    user\_interruption  & 166--193 & 0.385 & 0.427 & 0.461 & \textbf{0.076} \\
    user\_backchannel   &  78--92  & 0.269 & 0.326 & 0.315 & \textbf{0.057} \\
    talking\_to\_other  &  72--93  & 0.389 & 0.516 & 0.500 & \textbf{0.127} \\
    background\_speech  &  59--86  & 0.458 & 0.605 & 0.588 & \textbf{0.147} \\
    \bottomrule
  \end{tabular}
\end{table}

Table~\ref{tab:vad_variance} shows ATR for each backend. The variation across
the three systems is small relative to the ASR gap in
Table~\ref{tab:asr_sensitivity}. On \emph{user\_interruption}, where TOR is
most directly interpretable, the three VAD backends span a range of 7.6 pp ---
compared to a 23.5 pp gap between the two ASR systems on the same audio.
\emph{user\_backchannel} is similarly tight at 5.7 pp versus 23.5 pp.

The largest VAD disagreement is on \emph{background\_speech} (14.7~pp). This
is also the noisiest scenario --- ambient speech overlapping the model output
--- and VAD systems are known to differ more on non-speech noise than on clean
voiced speech. Even here, the VAD range remains smaller than the ASR gap on
\emph{talking\_to\_other} (22~pp) and \emph{user\_backchannel} (23.5~pp).

\begin{table}[t]
  \centering
  \caption{Maximum pairwise backend disagreement: ATR (three VAD systems)
           versus TOR (two ASR systems) on the same Moshi outputs.}
  \label{tab:comparison}
  \begin{tabular}{lcc}
    \toprule
    Scenario & TOR Max~$\Delta$ & ATR Max~$\Delta$ \\
    \midrule
    user\_interruption  & $0.000^{\dagger}$ & \textbf{0.076} \\
    user\_backchannel   & 0.235             & \textbf{0.057} \\
    talking\_to\_other  & 0.220             & \textbf{0.127} \\
    background\_speech  & 0.170             & \textbf{0.147} \\
    \bottomrule
  \end{tabular}
  {\small $^{\dagger}$Both ASR systems return TOR~$= 0$, but AssemblyAI
  produced empty transcripts for 54\% of samples in this scenario.}
\end{table}

Table~\ref{tab:comparison} puts the two metrics side by side. ATR is more
stable than TOR in three of four scenarios. The one apparent exception ---
\emph{user\_interruption} --- is where TOR happens to agree at zero, not
because the metric is stable but because Moshi genuinely does not yield the
floor on this task, and both ASR systems independently confirm it by
returning near-zero transcripts.

\subsection{Moshi ATR summary}

\begin{table}[t]
  \centering
  \caption{Moshi ATR (Silero VAD) on Full-Duplex-Bench v1.5 across all
           four overlap scenarios.}
  \label{tab:moshi_atr}
  \begin{tabular}{lcc}
    \toprule
    Scenario & ATR (Silero) & $n$ \\
    \midrule
    user\_interruption  & 0.385 & 166 \\
    user\_backchannel   & 0.269 &  78 \\
    talking\_to\_other  & 0.389 &  72 \\
    background\_speech  & 0.458 &  59 \\
    \bottomrule
  \end{tabular}
\end{table}

Moshi yields the floor in roughly 38.5\% of \emph{user\_interruption} scenarios
as measured by ATR, suggesting a partial floor-holding tendency consistent
with the behaviour reported for Moshi in~\cite{lin2025fullduplexv15}.
ATR values on the non-interruption scenarios (0.27--0.46) reflect cases where
the model continued or resumed speech after the overlap window, which is the
correct behaviour for backchannels and third-party speech --- though ATR alone
cannot distinguish between these cases without additional context. Pairing ATR
with the categorical behaviour labels from FDB v1.5 gives a more complete
picture.

%=============================================================================
\section{Discussion}
\label{sec:discussion}

\subsection{Why VAD is more stable than ASR here}

ASR systems make lexical decisions. A word at the tail of an utterance --- often
short, spoken at lower energy as the model winds down --- sits right at the
confidence boundary. One system transcribes it; another discards it. That
one-word difference shifts $t^{\text{asr}}_i$ by hundreds of milliseconds and
can flip the binary indicator in Equation~\ref{eq:tor}. Multiplied across a
dataset, these flips produce the gaps in Table~\ref{tab:asr_sensitivity}.

VAD systems do not have this problem because they do not make word-level
decisions. They classify frames as voiced or unvoiced based on acoustic
features, and the three systems we tested use quite different approaches ---
an LSTM, a GMM, and a transformer. That they agree as closely as they do
suggests the acoustic signal itself is informative for this task, and that
the disagreement in ASR-based TOR is mostly an artefact of the transcription
step rather than genuine ambiguity in the audio.

\subsection{ATR vs TOR: what they actually measure}

There is a real difference between asking ``when did the model's last word
end'' and ``when did the model's audio go silent.'' Trailing sounds ---
residual codec noise, unvoiced fricatives, breathing --- can push
$t^{\text{vad}}_i$ slightly later than $t^{\text{asr}}_i$. In most cases
this is harmless, but on clips where the overlap window is short, it could
flip a sample from ATR~$= 1$ to ATR~$= 0$.

Arguably, ATR is the more valid measure from a perceptual standpoint. A listener
perceives a model as having stopped when the audio falls silent, not when its
last transcribed word ended. Whether this aligns with human judgements of
natural floor-yielding is an empirical question. We are conducting a human
annotation study to investigate this, and will report results in a subsequent
version of this work.

\subsection{Limitations}

ATR cannot distinguish between a model that yields the floor intentionally
and one that simply produces no output. A model that fails to generate any
speech will score ATR~$= 1$ by default, which inflates its apparent
floor-yielding rate. Reporting a speech-presence rate alongside ATR addresses
this: if only 30\% of samples contain any detected speech, an ATR of 0.9 means
something very different than if 95\% do.

VAD systems are also not uniformly accurate across audio conditions. On noisy
scenarios like \emph{background\_speech}, we observed 14.7 pp disagreement
across backends --- still smaller than the ASR gap, but not negligible. Users
should select a VAD backend appropriate to their audio conditions and report
which one they used.

Finally, the VAD coverage varies across backends (Table~\ref{tab:vad_variance}
shows $n$ ranging from 59 to 193 depending on backend and scenario), because
systems differ in their sensitivity to low-energy or short speech segments.
This affects the denominator of ATR and should be reported alongside the rate.

%=============================================================================
\section{Conclusion}
\label{sec:conclusion}

TOR is a reasonable metric with a fragile implementation. Tying it to ASR
transcripts introduces variance that has nothing to do with how good a model
is at yielding the floor. Two standard ASR systems disagree by up to 23.5
percentage points on the same audio --- enough to change the story a paper
tells about a model's behaviour.

ATR keeps the same intuition as TOR and swaps the ASR step for VAD. Three
independent VAD backends agree within 7.6 pp on \emph{user\_interruption}
and within 14.7 pp across all scenarios. It runs locally, requires no
transcription, and works on any language. We are not arguing that VAD is
perfect, or that ATR captures everything TOR does. But for the specific
question of whether the model's audio fell silent before the overlap window
closed, VAD gives a more stable answer than ASR --- and stability is what
you need when comparing systems across labs and papers.

Code and evaluation scripts are available at:
\url{https://github.com/ifeanyidike/full-duplex-bench-atr}

%=============================================================================
\section{Acknowledgements}

We thank Guan-Ting Lin for making Full-Duplex-Bench publicly available.
% To be completed.

%=============================================================================
\bibliographystyle{IEEEtran}
\bibliography{atr_refs}

\end{document}
